\par\medskip
且对每个非空子集 $B\subset S$, 有

\begin{equation}\tag{4.5.20}
\bigoplus_{A\supset B}L^A(b)
=L(b)\cap\left(\prod_{i\in B}X_i\right)K[[X]].
\end{equation}

显然每个子空间 $\bigoplus_{A\supset B}L^A(b)$ 都在 $\alpha$ 下稳定.
因此 $\bigoplus_{A\not\supset B}L^A(b)$ 也如此, 因为它是前一类型的子空间之和.
定义

\begin{equation}\tag{4.5.21}
\alpha_B=\text{$\alpha$ 在下列空间上诱导的 $\tau^{-1}$ 线性自同态}
\quad
\frac{\displaystyle\bigoplus_{A\supset B}
L^A\!\left(\frac1{p-1}\right)}
{\displaystyle\bigoplus_{\substack{A\supset B\\A\ne B}}
L^A\!\left(\frac1{p-1}\right)}
\simeq L^B\!\left(\frac1{p-1}\right).
\end{equation}

明确地, $\alpha_B$ 由下式给出

\begin{equation}\tag{4.5.22}
\alpha_B=\frac{\tau^{-1}}p\,\psi_p\,H_B(X),
\end{equation}

$L^B$ 在此运算下稳定, 其中 $H_B(X)$ 由 $H(X)$ 中令 $i\notin\{0\}\cup B$ 时
$X_i=0$ 得到.

再次改写 (4.5.16) 得

\begin{equation}\tag{4.5.23}
\operatorname{card}V^*(\mathbb F_{q^\nu})
=(q^\nu-1)^{n+1}
+(q^\nu-1)^{n+2}
\sum_{\substack{A\subset S\\A\ne\varnothing}}
\operatorname{tr}\!\left((\alpha_A)^{\operatorname{ord}(q^\nu)}\right).
\end{equation}

对 $A\subset S$, $A\ne\varnothing$, 定义超曲面

\begin{equation}\tag{4.5.24}
V_A\subset\mathbb P^{\operatorname{card}(A)-1}
\simeq
\mathbb P^{n+1}\cap
\{X_i=0\text{ 对每个 }i\notin A\}
\end{equation}

其中多项式 $f_A(X)$ 由 $f(X)$ 中令 $i\notin A$ 时 $X_i=0$ 得到.

由从 $f(X)$ 形成级数 $H$ 与将部分变量置为 $0$ 的相容性,
(4.5.23) 对每个 $A\subset S$, $A\ne\varnothing$ 给出

\begin{equation}\tag{4.5.25}
\operatorname{card}V_A^*(\mathbb F_{q^\nu})
=(q^\nu-1)^{\operatorname{card}(A)-1}
+(q^\nu-1)^{\operatorname{card}(A)}
\sum_{B\subset A}
\operatorname{tr}\!\left((\alpha_B)^{\operatorname{ord}(q^\nu)}\right).
\end{equation}

另一方面,

\begin{equation}\tag{4.5.26}
\operatorname{card}V(\mathbb F_{q^\nu})
=\sum_{\substack{A\subset S\\A\ne\varnothing}}
\operatorname{card}V_A^*(\mathbb F_{q^\nu}).
\end{equation}

因此经简单求和, (4.5.25) 给出

\begin{equation}\tag{4.5.27}
\begin{aligned}
\operatorname{card}V(\mathbb F_{q^\nu})
={}&1+q^\nu+q^{2\nu}+\cdots
+q^{\nu(\operatorname{card}(S)-1)}\\
&+\sum_{\substack{B\subset S\\B\ne\varnothing}}
(q^\nu-1)^{\operatorname{card}(B)}
q^{\nu\operatorname{card}(S-B)}
\operatorname{tr}\!\left((\alpha_B)^{\operatorname{ord}(q^\nu)}\right).
\end{aligned}
\end{equation}

由于 $\alpha$ 从而各 $\alpha_B$ 都是 $\tau^{-1}$ 线性的,
算子 $\alpha_B^{\operatorname{ord}(q)}$ 为 $K$ 线性的. 定义特征级数

\[
\Delta_B(t)\in1+tK[[t]]
\]

其定义式为

\begin{equation}\tag{4.5.28}
\Delta_B(t)=\det\!\left(
1-t\alpha_B^{\operatorname{ord}(q)}
\mathrel{\Big|}L^B\!\left(\frac1{p-1}\right)\right).
\end{equation}

引入 $1+tK[[t]]$ 的下述自同态:

\begin{equation}\tag{4.5.29}
\varphi:g(t)\longmapsto g^\varphi(t)=g(qt)
\end{equation}

以及

\begin{equation}\tag{4.5.30}
\delta=1-\varphi:g(t)\longmapsto
g^\delta(t)=\frac{g(t)}{g(qt)}.
\end{equation}
